\section{Related Work}\label{sec-3}

\subsection{Rare Disease LLM Benchmarks}\label{sec-3-1}

\textbf{Existing benchmarks are LLM-centric, not agent-native.} RareBench
\citep{rarebench2024} establishes a multi-subset rare-disease diagnostic
protocol with rank-based metrics. RareArena \citep{rarearena2025} scales to
free-text case reports, while Phenopacket-Store and subsequent Phenopacket
studies provide structured HPO-based evaluation
\citep{phenopacketstore2025,reese2026,chimirri2025}. MIMIC-RD
\citep{mimicrd2026} introduces a note-based MIMIC-IV rare-disease task; unlike
that resource, our local MIMIC probe has no notes. These benchmarks primarily
evaluate static input-to-output diagnosis and do not jointly expose interactive
tool use, cost/latency, reasoning-trace evaluation, and pass-\(k\) reliability.
A systematic review of 15 studies (19 system--dataset entries) found all 19 at
high risk of bias and highlighted substantial benchmark-dependent
heterogeneity \citep{sysreview2026}.

\subsection{Agent Benchmarks in Other Domains}\label{sec-3-2}

\textbf{Agent benchmarks elsewhere have matured beyond static evaluation; rare-disease benchmarks have not.} \(\tau\)-bench formalizes repeated-trial
reliability for tool-using agents \citep{taubench2024}; AgentBoard introduces
partial-credit progress metrics \citep{agentboard2024}; and SWE-bench makes
end-to-end issue resolution the unit of evaluation \citep{swebench2024}. In
medicine, MedAgentBench covers tool-augmented FHIR tasks
\citep{medagentbench2025}, while MedHELM extends scenario-by-metric evaluation
\citep{medhelm2025}. We adopt three corresponding patterns: process-aware
evaluation, pass-\(k\) reliability, and explicit cost/latency accounting.

\subsection{Rare-Disease and Medical Agent Systems}\label{sec-3-3}

\textbf{Eight systems are relevant to rare-disease diagnosis or transferable to it, but they do not share an evaluation matrix.} DeepRare integrates more
than 40 tools and evaluates heterogeneous text, HPO, and genetic inputs
\citep{deeprare2026}. MAI-DxO coordinates a multi-role diagnostic panel
\citep{maidxo2025}; RDMA specializes in phenotype extraction
\citep{rdma2025}; and VC-RDAgent provides an offline ontology-based path
\citep{vcrdagent}. From general medicine, MDAgents adapts collaboration mode
\citep{mdagents2024}, MedAgents orchestrates role-playing experts
\citep{medagents2024}, and AgentClinic evaluates simulated clinical dialogue
\citep{agentclinic2024}. LIRICAL supplies a classical likelihood baseline
\citep{lirical2020}. Each was originally evaluated under a different input,
dataset, or metric, leaving cross-system claims unverifiable on common ground.

We position this work as filling that exact gap.

